\documentclass[12pt,a4paper,openany]{report}
\usepackage{tgcommon}

\title{Estimação Neural de Processos de Hawkes em Dados de Criptomoedas de Alta Frequência}
\author{João Matheus Del Vecchio França Barbosa}
\date{2026}

\begin{document}
\tgtitlepage
  {João Matheus Del Vecchio França Barbosa}
  {Estimação Neural de Processos de Hawkes em Dados de Criptomoedas de Alta Frequência}
  {Rascunho de dissertação --- reprodução de LSTM e extensão com atenção}
  {2026}

\pagenumbering{roman}
\chapter*{Resumo}

Processos de Hawkes descrevem sequências de eventos em que cada ocorrência eleva
temporariamente a probabilidade de novas ocorrências; em microestrutura financeira,
representam o agrupamento de alterações de preço. Este trabalho reproduz e estende a
estimação amortizada, por redes neurais, dos parâmetros de um Hawkes bidimensional
simétrico proposta por \citet{lee2023}, com dados de alta frequência dos futuros
perpétuos BTCUSDT e ETHUSDT da Binance. Foi construída uma pipeline reproduzível que
reconstrói o mid-price em grade de 100~ms (1.826.904 eventos na janela principal e
1.288.352 na de robustez), simula trajetórias por \emph{thinning} de Ogata, estima
parâmetros por máxima verossimilhança (MLE) e treina uma LSTM, como no artigo, e um
encoder Transformer. No protocolo sintético congelado (75.000/10.000/15.000
trajetórias de 1.000 eventos e cinco seeds por rede), a LSTM obteve MSE macro de
0,0196 e o Transformer, 0,0237; com cinco seeds, a diferença não é significativa
($p=0{,}19$), e o ensemble da LSTM supera o do Transformer em 70\% das trajetórias.
O MSE macro do MLE (0,506) é dominado por 1\% das trajetórias, nas quais $\beta$ é
pouco identificado por causa da excitação fraca; na trajetória típica, o MLE é tão
preciso quanto as redes e melhor em $\alpha_{\mathrm{self}}$ e
$\alpha_{\mathrm{cross}}$. No ponto fixo publicado por \citet{lee2023}, o
experimento reproduz o padrão do artigo: redes com viés de encolhimento e MLE quase
não viesado. Um benchmark controlado mostrou custo de cerca de 10~ms por trajetória
para o MLE e de 5 a 10~ms para as redes em lote unitário; a vantagem de velocidade
das redes aparece apenas na inferência em lote. Por fim, o MLE em 1.063 janelas reais
de calibração indicou excitação cruzada praticamente nula e $\beta$ mediano de 2,7,
fora do suporte usado no treino; por isso, a adaptação empírica exige recalibrar o
prior ou treinar com rótulos MLE em janelas reais. Conclui-se que a atenção não
superou a LSTM nesse problema e que a principal vantagem da estimação amortizada é a
robustez e o custo em lote, não a precisão por trajetória.

\noindent\textbf{Palavras-chave:} processos de Hawkes; alta frequência; criptomoedas;
LSTM; Transformer; estimação amortizada.

\chapter*{Abstract}

Hawkes processes describe event sequences in which each occurrence temporarily raises
the probability of further occurrences; in market microstructure they represent the
clustering of price changes. This work reproduces and extends the amortized neural
estimation of a symmetric bivariate Hawkes process proposed by \citet{lee2023}, using
high-frequency data from Binance BTCUSDT and ETHUSDT perpetual futures. A reproducible
pipeline reconstructs the mid-price on a 100~ms grid (1,826,904 events in the primary
window and 1,288,352 in the robustness window), simulates paths by Ogata thinning,
estimates parameters by maximum likelihood (MLE), and trains the paper's LSTM and a
Transformer encoder. Under the frozen synthetic protocol (75,000/10,000/15,000 paths
of 1,000 events, five seeds per network), the LSTM reached a macro MSE of 0.0196 and
the Transformer 0.0237; with five seeds the difference is not significant ($p=0.19$),
and the LSTM ensemble beats the Transformer ensemble on 70\% of the paths. The MLE
macro MSE (0.506) is driven by 1\% of the paths, in which $\beta$ is weakly identified
because excitation is low; on a typical path the MLE is as accurate as the networks
and better for $\alpha_{\mathrm{self}}$ and $\alpha_{\mathrm{cross}}$. At the fixed
point published by \citet{lee2023}, the experiment reproduces the paper's pattern:
shrinkage bias for the networks and a nearly unbiased MLE. A controlled benchmark
measured about 10~ms per path for the MLE and 5 to 10~ms for the networks at batch
size one, so the speed advantage of the networks appears only in batched inference.
Finally, MLE fits on 1,063 real calibration windows show near-zero cross-excitation
and a median $\beta$ of 2.7, outside the training support; the empirical adaptation
therefore requires a recalibrated prior or training on MLE-labelled real windows.
Attention did not outperform the LSTM on this problem, and the main benefit of
amortized estimation is robustness and batched cost rather than per-path accuracy.

\noindent\textbf{Keywords:} Hawkes processes; high frequency; cryptocurrencies; LSTM;
Transformer; amortized inference.

\tableofcontents
\listoftables
\listoffigures
\clearpage
\pagenumbering{arabic}

\chapter{Introdução}

\section{Contexto}

Preços financeiros não mudam em intervalos regulares. Em escalas curtas, observa-se
uma sequência de eventos: alteração para cima, alteração para baixo, negociação,
inclusão ou remoção de ordem. A atividade alterna períodos calmos e rajadas. Essa
dependência temporal é chamada de \emph{clustering} e está ligada à volatilidade
\citep{bouchaud2018}.

Um processo de Poisson homogêneo supõe taxa constante e incrementos independentes.
Essa hipótese é útil como ponto de partida, mas incapaz de representar o fato de que
uma mudança recente costuma ser seguida por mais mudanças. O processo de Hawkes
\citep{hawkes1971,hawkes1971b} corrige essa limitação ao tornar a intensidade
dependente da própria história e tornou-se um modelo de referência em finanças de
alta frequência \citep{bacry2015,hawkes2018}.

O desafio prático é estimar os parâmetros. MLE possui base estatística clara, porém
exige otimização iterativa para cada nova janela. \citet{lee2023} propõe pagar um custo
de treinamento uma vez e usar uma LSTM para mapear diretamente a sequência observada
nos parâmetros. Essa estratégia é \emph{inferência amortizada}: várias inferências
futuras reutilizam o conhecimento aprendido \citep{cranmer2020}.

\section{Problema}

O artigo de referência utiliza ações dos Estados Unidos entre 2018 e 2019. Seus dados
NBBO não foram publicados, e detalhes importantes do protocolo estão ausentes. Além
disso, somente uma arquitetura recorrente compacta foi estudada. Criptoativos oferecem
um caso de teste natural: negociação contínua, 24 horas por dia, e dados públicos de
melhor oferta \citep{makarov2020}. O presente trabalho pergunta:

\begin{quote}
Uma arquitetura com atenção consegue estimar parâmetros de um Hawkes 2D em movimentos
de preço de criptomoedas com acurácia e estabilidade iguais ou superiores à LSTM de
Lee, preservando inferência mais rápida que MLE?
\end{quote}

\section{Hipóteses}

As hipóteses foram registradas antes da execução final e são mantidas como foram
escritas:

\begin{description}
  \item[H1.] Dados públicos USD-M da Binance permitem construir sequências equivalentes
  em estrutura às entradas de Lee.
  \item[H2.] A LSTM alcança erro da mesma ordem do MLE com latência substancialmente menor.
  \item[H3.] A atenção melhora erro ou robustez em sequências com dependências longas.
\end{description}

H3 pode ser rejeitada. Um resultado negativo, se obtido sob protocolo congelado, é
cientificamente informativo.

\section{Objetivos}

O objetivo geral é reproduzir e estender a estimação neural de Hawkes em cripto.
Objetivos específicos:
\begin{enumerate}
  \item construir e auditar o dataset de top-of-book;
  \item implementar Hawkes 2D, simulação e MLE;
  \item reproduzir a LSTM do artigo;
  \item adaptar a pipeline a BTCUSDT e ETHUSDT;
  \item projetar Transformer comparável;
  \item avaliar acurácia, validade, custo e robustez;
  \item aplicar o melhor estimador à volatilidade Hawkes.
\end{enumerate}
Os objetivos 4 e 7 dependem da etapa empírica descrita no
Capítulo~\ref{cap:futuro}.

\section{Contribuições}

A primeira contribuição é de reprodução: código aberto, registro das lacunas do
artigo e reprodução quantitativa da distribuição amostral no ponto fixo publicado. A
segunda é metodológica: separação explícita entre reprodução sintética, adaptação
empírica e extensão, com testes estatísticos pareados e métricas robustas que mostram
por que o MSE agregado favorece as redes. A terceira é arquitetural: comparação justa
de LSTM e Transformer, com resultado negativo para a atenção. A quarta é empírica:
caracterização dos eventos de mid-price de BTCUSDT e ETHUSDT e evidência de que seus
parâmetros estimados ficam fora da região típica de ações.

\chapter{Fundamentação teórica}

\section{Processos pontuais}

Se \(t_1<t_2<\cdots\) são instantes de evento, o processo de contagem é
\[
  N(t)=\sum_k\mathbf{1}_{\{t_k\le t\}}.
\]
A intensidade condicional \(\lambda(t\mid\mathcal H_t)\) descreve a taxa instantânea
esperada dado o histórico \(\mathcal H_t\) \citep{daley2003}. Informalmente,
\[
  \Pr[N(t+\mathrm dt)-N(t)=1\mid\mathcal H_t]
  \approx \lambda(t)\,\mathrm dt.
\]
Não se trata de probabilidade em unidade arbitrária: intensidade possui dimensão de
eventos por unidade de tempo.

\section{Poisson como modelo nulo}

No Poisson homogêneo, \(\lambda(t)=\lambda_0\). Tempos entre eventos são exponenciais
independentes. O modelo não produz rajadas além do acaso. Portanto, é um controle útil:
se Hawkes não superar Poisson em diagnóstico, a autoexcitação não está justificada.

\section{Hawkes univariado}

Com kernel exponencial,
\[
  \lambda(t)=\mu+\sum_{t_k<t}\alpha e^{-\beta(t-t_k)}.
\]
\(\mu\) é atividade exógena; \(\alpha\) é o salto imediato; \(\beta\) controla a
velocidade de esquecimento. A integral do kernel é \(\alpha/\beta\), a razão de
ramificação \citep{laub2021}.

\section{Hawkes bidimensional simétrico}

Altas e baixas são processos \(N_1,N_2\):
\begin{align}
\lambda_1(t)&=\mu+\sum_{t_k^1<t}\alpha_{\mathrm{self}}e^{-\beta(t-t_k^1)}
 +\sum_{t_k^2<t}\alpha_{\mathrm{cross}}e^{-\beta(t-t_k^2)},\\
\lambda_2(t)&=\mu+\sum_{t_k^2<t}\alpha_{\mathrm{self}}e^{-\beta(t-t_k^2)}
 +\sum_{t_k^1<t}\alpha_{\mathrm{cross}}e^{-\beta(t-t_k^1)}.
\end{align}
O modelo reduz uma matriz geral a quatro parâmetros. Essa parcimônia facilita
simulação e comparação com \citet{lee2023}. Em microestrutura, excitação cruzada
dominante indica reversão (uma alta seguida de baixa), e autoexcitação dominante
indica persistência de direção \citep{bacry2013}.

\section{Estabilidade}

A matriz de massa dos kernels é
\[
 \mathbf G=\frac{1}{\beta}
 \begin{bmatrix}
  \alpha_{\mathrm{self}} & \alpha_{\mathrm{cross}}\\
  \alpha_{\mathrm{cross}} & \alpha_{\mathrm{self}}
 \end{bmatrix}.
\]
Seu maior autovalor é
\[
 g=\frac{\alpha_{\mathrm{self}}+\alpha_{\mathrm{cross}}}{\beta}.
\]
Exige-se \(g<1\). Perto de um, cascatas duram mais; acima de um, não há regime
estacionário usual \citep{hardiman2013,filimonov2012}.

\section{Marcas e preço}

Um evento tem tempo, tipo e marca. A marca é o tamanho da mudança em ticks. O artigo
preserva a marca para volatilidade, mas supõe que a intensidade não depende dela. O
preço filtrado pode ser escrito, em ticks, como diferença acumulada entre altas e baixas.

\section{Volatilidade}

Em um modelo mínimo, se o preço se move por um tick \(\vartheta\) a uma taxa \(\phi\),
\[
 \sigma^2\approx\phi\vartheta^2.
\]
Assim, modelar intensidade é modelar atividade e, por consequência, uma componente da
volatilidade. A expressão de Hawkes marcada requer momentos das marcas e parâmetros
estimados \citep{leeseo2017,lee2024}. A realized volatility será o benchmark
\citep{andersen2003,barndorff2002}; em alta frequência, ruído de microestrutura exige
cuidado na escolha da amostragem \citep{zhang2005}.

\section{Máxima verossimilhança}

Para horizonte \(T\),
\[
 \ell(\theta)=\sum_k\log\lambda_{c_k}(t_k)
 -\sum_i\int_0^T\lambda_i(u)\,du.
\]
O primeiro termo recompensa intensidade alta nos eventos; o segundo penaliza intensidade
alta em todo o intervalo. Maximizar apenas o primeiro faria intensidades tenderem ao
infinito. O compensador produz equilíbrio. Para kernel exponencial, a soma sobre o
passado obedece a uma recursão de primeira ordem, e \(\ell\) e seu gradiente podem ser
calculados em \(O(n)\) operações \citep{ozaki1979}. Sob condições de regularidade, o
MLE é consistente e assintoticamente eficiente \citep{ogata1978}.

\section{Identificabilidade em amostras finitas}
\label{sec:identificabilidade}

As propriedades assintóticas não garantem bom comportamento com 1.000 eventos. Quando
a excitação é fraca (\(\alpha_{\mathrm{self}}+\alpha_{\mathrm{cross}}\) pequeno), o
kernel quase não altera a intensidade, e a verossimilhança fica plana na direção de
\(\beta\): valores muito diferentes de \(\beta\) explicam os dados quase igualmente bem.
No limite \(\alpha\to0\), \(\beta\) deixa de ser identificado. Perto da criticidade,
a separação entre \(\mu\) e \(g\) também se deteriora, pois a taxa média
\(2\mu/(1-g)\) é bem determinada, mas cada fator isoladamente não é. Esses efeitos
reaparecem nos resultados do Capítulo~\ref{cap:resultados}.

\section{Estimação amortizada como estimador de Bayes}
\label{sec:bayes}

Na abordagem de \citet{lee2023}, uma rede \(f_\phi\) é treinada com pares
\((\theta_i,x_i)\), em que \(\theta_i\) é sorteado de uma distribuição de treino
\(\pi\) e \(x_i\) é simulado do modelo com parâmetro \(\theta_i\). Minimizar o erro
quadrático médio,
\[
  \min_\phi\;\E_{\theta\sim\pi}\,\E_{x\mid\theta}
  \bigl\lVert f_\phi(x)-\theta\bigr\rVert^2,
\]
é minimizar o risco de Bayes sob perda quadrática, cuja solução é a média a
posteriori \(\E_\pi[\theta\mid x]\) \citep{lehmann1998}. Logo, uma rede bem treinada
aproxima um estimador de Bayes sob o prior \(\pi\)
\citep{sainsburydale2024,lenzi2023}. Duas consequências orientam este trabalho.
Primeiro, no teste sorteado do mesmo \(\pi\), nenhum estimador tem menor MSE médio
que a média a posteriori, de modo que superar o MLE nesse teste é esperado e não prova
superioridade em sentido frequentista. Segundo, a rede não extrapola: suas estimativas
encolhem para a região de alta densidade de \(\pi\), e dados reais fora do suporte do
prior produzem estimativas viesadas.

\section{LSTM}

Redes recorrentes atualizam um estado \(h_t=f(x_t,h_{t-1})\). LSTM introduz portas de
entrada, esquecimento e saída para controlar fluxo de informação e reduzir o problema
de gradientes \citep{hochreiter1997}. Em Lee, a sequência contém interchegada e tipo;
o estado final alimenta a regressão dos quatro parâmetros. O artigo também testou GRU
\citep{cho2014}, com resultado ligeiramente inferior.

\section{Atenção e Transformer}

Atenção calcula combinações ponderadas:
\[
 \operatorname{Attention}(Q,K,V)=
 \operatorname{softmax}\left(\frac{QK^\top}{\sqrt{d_k}}\right)V.
\]
Cada posição consulta todas as outras. Transformer combina múltiplas cabeças,
feed-forward, resíduos e normalização \citep{vaswani2017}. A vantagem potencial é
capturar dependências distantes; o custo potencial é memória \(O(n^2)\). A posição da
normalização em relação aos resíduos afeta a estabilidade do treino: a variante
\emph{post-LN}, usada aqui, costuma exigir aquecimento da taxa de aprendizado
\citep{xiong2020}.

\chapter{Trabalhos relacionados}

\section{Processos de Hawkes em finanças}

Depois da formulação original \citep{hawkes1971,hawkes1971b}, modelos auto e
mutuamente excitantes passaram a descrever eventos de mercado em tempo contínuo
\citep{bowsher2007,embrechts2011}, contágio entre mercados \citep{aitsahalia2015} e
endogeneidade, medida pela razão de ramificação \citep{filimonov2012,hardiman2013}.
\citet{bacry2013} usam um Hawkes bidimensional de altas e baixas para reproduzir o
ruído de microestrutura do preço; a excitação cruzada captura a reversão em escalas
curtas. \citet{bacry2015} e \citet{hawkes2018} revisam esse campo, e
\citet{mutoke2011} oferece uma introdução didática. \citet{leeseo2017} conectam o
Hawkes marcado ao preço em ticks e derivam a volatilidade de Hawkes, aplicada a preços
filtrados por \citet{lee2024}. No Brasil, \citet{primo2018} analisa o agrupamento
intradiário de ações com MLE, e \citet{silva2026} sistematiza inferência, diagnósticos
e aplicações financeiras.

\section{Estimação de processos de Hawkes}

O MLE para kernels exponenciais remonta a \citet{ozaki1979}, com propriedades
assintóticas estabelecidas por \citet{ogata1978}. A simulação por \emph{thinning}
\citep{ogata1981} e o diagnóstico por resíduos reescalados \citep{ogata1988,brown2002}
completam a caixa de ferramentas clássica \citep{laub2021}. A otimização é sensível à
inicialização e à geometria da verossimilhança, que pode ser quase plana em algumas
direções quando a série é curta (Seção~\ref{sec:identificabilidade}).

\section{Inferência amortizada e baseada em simulação}

Quando a verossimilhança é cara ou intratável, redes treinadas em simulações
substituem a otimização por uma passagem direta \citep{cranmer2020}. Redes treinadas
em simulações já estimam parâmetros de modelos espaciais e de valores extremos
\citep{gerber2021,lenzi2023}; com perda quadrática, elas aproximam estimadores de Bayes
\citep{sainsburydale2024}, e métodos como BayesFlow aprendem a posterior inteira
\citep{radev2022}. \citet{lee2023} aplica a ideia ao Hawkes 2D com
uma LSTM. \citet{lambe2025} usam redes para o Hawkes observado apenas por contagens
em intervalos, caso em que a verossimilhança exata é intratável e o MLE deixa de ser
uma alternativa direta.

\section{Processos pontuais neurais}

Outra linha usa redes para modelar a própria intensidade, sem a forma paramétrica de
Hawkes: RMTPP \citep{du2016}, o Neural Hawkes Process \citep{mei2017} e variantes com
atenção, como o Transformer Hawkes Process \citep{zuo2020} e o Self-Attentive Hawkes
Process \citep{zhang2020}; \citet{shchur2021} revisam a área. O objetivo é diferente
do deste trabalho: esses modelos maximizam verossimilhança preditiva, enquanto aqui a
rede estima os quatro parâmetros interpretáveis de um modelo clássico.

\section{Hawkes em criptoativos}

Com retornos de um minuto de 17 criptomoedas negociadas na Binance,
\citet{taddeo2024} mostram que o Hawkes descreve melhor que o Poisson a ocorrência de
saltos nos log-retornos. O mercado de cripto também apresenta
segmentação e arbitragem limitada entre bolsas \citep{makarov2020}, e os futuros
perpétuos, sem vencimento e com taxa de financiamento, concentram grande parte da
liquidez \citep{he2022perpetual}.

\section{Posicionamento}

\begin{table}[htbp]
\centering
\caption{Posicionamento do trabalho em relação à literatura próxima.}
\label{tab:posicionamento}
\small
\begin{tabularx}{\linewidth}{lllX}
\toprule
Trabalho & Estimação & Arquitetura & Mercado\\
\midrule
\citet{bacry2013} & MLE & -- & futuros (Bund, DAX)\\
\citet{leeseo2017} & MLE & -- & ações\\
\citet{primo2018} & MLE & -- & ações (Brasil)\\
\citet{taddeo2024} & MLE & -- & cripto (saltos em retornos de 1 min)\\
\citet{lee2023} & rede & LSTM & ações (EUA)\\
\citet{lambe2025} & rede & MLP & sintético (contagens)\\
\citet{zuo2020} & intensidade neural & Transformer & genérico\\
Este trabalho & rede e MLE & LSTM e Transformer & cripto (futuros perpétuos)\\
\bottomrule
\end{tabularx}
\end{table}

O presente trabalho combina reprodução controlada, dados públicos de cripto e
comparação com atenção, e acrescenta à avaliação de \citet{lee2023} métricas robustas,
testes pareados e a leitura da rede como estimador de Bayes.

\chapter{Dados}

\section{Fonte e mercado}

Foram avaliados spot e futuros. A necessidade de melhor bid/ask eliminou o acervo spot
oficial, que não possui \dataset{bookTicker} histórico equivalente. Futuros USD-M
fornecem \dataset{bookTicker} e \dataset{aggTrades}. Os símbolos são BTCUSDT e
ETHUSDT \citep{binancePublicData}.

\section{Schema}

\dataset{bookTicker} inclui identificador, melhores preços e quantidades de compra e
venda, tempo de transação e tempo de evento. O schema canônico possui timestamp UTC,
bid, ask e quantidades. Auditorias verificam preço positivo, bid menor ou igual a ask,
ordenação e duplicidade temporal.

\section{Amostragem}

Atualizações podem ocorrer a cada milissegundo, muito acima da escala escolhida.
Agrupam-se timestamps em bins de 100 ms e mantém-se o último estado. Esse filtro segue
o protocolo de Lee e reduz ruído de ultra-alta frequência.

\section{Datasets produzidos}

\begin{table}[htbp]
\centering
\caption{Cobertura dos datasets.}
\label{tab:cobertura}
\begin{tabular}{lrrr}
\toprule
Janela & Dias-símbolo & Eventos & Segmentos de 1.000\\
\midrule
Principal (jan. 2024) & 28 & 1.826.904 & 1.813\\
Robustez (out. 2023) & 14 & 1.288.352 & 1.281\\
\bottomrule
\end{tabular}
\end{table}

\section{Características dos eventos}
\label{sec:caracteristicas}

A Tabela~\ref{tab:caracteristicas} resume os eventos do bloco de calibração (15 a
22 de janeiro de 2024). Três fatos distinguem esses dados do cenário de
\citet{lee2023}. Primeiro, cerca de 76\% das mudanças repetem a direção da anterior,
o que indica persistência, e não reversão. Segundo, a mediana do salto é de 10 a 15
ticks e só 10--13\% dos saltos têm um tick: na escala de 100~ms, esses contratos se
comportam como ativos de tick pequeno. Terceiro, 37--39\% das interchegadas são
exatamente 0,1~s, o passo da grade; a amostragem discretiza fortemente os tempos,
enquanto o modelo e as simulações supõem tempo contínuo.

\begin{table}[htbp]
\centering
\caption{Características dos eventos de mid-price, 15 a 22 de janeiro de 2024.}
\label{tab:caracteristicas}
\small
\begin{tabular}{lrr}
\toprule
 & BTCUSDT & ETHUSDT\\
\midrule
Eventos & 529.278 & 540.164\\
Fração de altas & 49,4\% & 49,5\%\\
Probabilidade de repetir a direção anterior & 76,0\% & 75,7\%\\
Interchegada mediana & 0,3 s & 0,3 s\\
Interchegadas iguais a 0,1 s (grade) & 39,5\% & 37,3\%\\
Salto mediano / percentil 90 (ticks) & 15,5 / 48 & 10 / 29\\
Saltos de um tick & 9,7\% & 12,6\%\\
\bottomrule
\end{tabular}
\end{table}

\chapter{Metodologia}

\section{Princípios}

O protocolo separa:
\begin{description}
  \item[R1:] informação explicitamente publicada;
  \item[R2:] reconstrução para preencher lacuna;
  \item[A1:] adaptação necessária para cripto;
  \item[E1:] extensão com Transformer.
\end{description}

\section{Pipeline}

\begin{enumerate}
  \item construir URL diária e baixar ZIP;
  \item calcular SHA-256 e registrar manifesto;
  \item extrair temporariamente o CSV;
  \item processar chunks de um milhão de linhas;
  \item reduzir atualizações à grade de 100 ms;
  \item produzir Parquet particionado por símbolo e data;
  \item apagar CSV temporário;
  \item criar janelas sem atravessar splits;
  \item ajustar MLE e calibrar faixas apenas no treino;
  \item gerar trajetórias sintéticas e treinar estimadores;
  \item selecionar hiperparâmetros na validação;
  \item abrir teste uma única vez.
\end{enumerate}

\section{Simulação por thinning}

Com estado atual, calcula-se um limite para a soma das intensidades. Sorteia-se espera
exponencial; a excitação decai; aceita-se o candidato proporcionalmente à nova
intensidade \citep{ogata1981}. O tipo é sorteado proporcionalmente a
\(\lambda_1,\lambda_2\). Após o evento, adiciona-se a coluna correspondente da matriz
\(\alpha\). Cada trajetória começa com histórico vazio, a mesma hipótese da
verossimilhança.

\section{MLE restrito}

Parâmetros positivos são obtidos por softplus. Em vez de otimizar \(\beta\) livremente,
usa-se margem positiva \(d\):
\[
(\mu,\alpha_s,\alpha_c,d)=\softplus(z),\qquad
\beta=\alpha_s+\alpha_c+d.
\]
L-BFGS-B é iniciado em cinco pontos: a taxa empírica de eventos e quatro perturbações
aleatórias no espaço irrestrito. O MLE nunca recebe o parâmetro verdadeiro como
inicialização. O gradiente é analítico, e as recorrências exponenciais são resolvidas
como sistemas bidiagonais. Registram-se convergência, iterações,
log-verossimilhança e tempo.

\section{Protocolo neural}

O alvo é \((\mu,\alpha_s,\alpha_c,\beta)\), e a perda é o MSE na escala original. A
entrada de cada evento é \((\log(1+\Delta t),\,s)\), com \(s=+1\) para alta e
\(s=-1\) para baixa; Lee usa interchegadas e tipos 1/2, de modo que essa transformação
é uma reconstrução (R2). As duas redes compartilham a cabeça de saída: softplus produz
\(\mu,\alpha_s,\alpha_c,d>0\) e \(\beta=\alpha_s+\alpha_c+d\), o que garante
\(g<1\) por construção, como na versão robusta de Lee.

O experimento sintético congelado contém 75.000 sequências de treino, 10.000 de
validação e 15.000 de teste, exclusivamente com 1.000 eventos. Os parâmetros são
sorteados com \(\mu\in[0{,}2;\,0{,}6]\), \(\beta\in[0{,}8;\,2{,}0]\), fração
autoexcitada \(\alpha_s/(\alpha_s+\alpha_c)\in[0{,}2;\,0{,}8]\) e \(g\) balanceado em
três estratos: baixo \([0{,}10;\,0{,}35)\), intermediário \([0{,}35;\,0{,}70)\) e quase
crítico \([0{,}70;\,0{,}95]\). A Tabela~\ref{tab:treino} resume o treinamento.

\begin{table}[htbp]
\centering
\caption{Configuração e custo do treinamento (cinco seeds por rede, Mac Apple
Silicon com MPS).}
\label{tab:treino}
\small
\begin{tabular}{lrr}
\toprule
 & LSTM (R1/R2) & Transformer (E1)\\
\midrule
Arquitetura & 2 LSTM $\times$ 12 + densa & encoder, 2 camadas, $d=32$\\
Cabeças / feed-forward & -- & 4 / 128\\
Parâmetros treináveis & 2.068 & 25.700\\
Otimizador & Adam, $10^{-3}$ & Adam, $10^{-3}$\\
Batch efetivo & 64 & 8 (físico 4 com acumulação)\\
Épocas máximas / paciência & 300 / 30 & 300 / 30\\
Melhor época (por seed) & 165, 164, 193, 239, 299 & 73, 42, 38, 60, 44\\
Tempo de treino por seed & 5,4--14,8 h & 16,6--28,8 h\\
Tempo total (5 seeds) & 41 h & 110 h\\
\bottomrule
\end{tabular}
\end{table}

O fluxo tensorial do Transformer parte de
\([B,L,2]\), com \(L=1000\). Uma projeção linear produz
\([B,L,32]\), à qual se soma o encoding posicional. Em cada camada, as quatro cabeças
operam em subespaços de dimensão oito; o bloco feed-forward expande 32 para 128 e
retorna a 32. A média temporal converte \([B,L,32]\) em \([B,32]\). O custo de memória
da atenção cresce quadraticamente com \(L\), razão pela qual seu batch é menor que o
da LSTM. Essa diferença de batch, sem busca de hiperparâmetros para nenhuma das redes,
é uma limitação da comparação.

\section{Ponto fixo do artigo (E2)}

Para comparar diretamente com a tabela de distribuição amostral de \citet{lee2023},
foram simuladas 2.000 trajetórias de 1.000 eventos no ponto publicado
\((\mu,\alpha_s,\alpha_c,\beta)=(0{,}3;\,0{,}4;\,0{,}7;\,1{,}5)\), que está dentro do
prior de treino. As redes foram aplicadas sem retreino, e o MLE foi ajustado em duas
versões: a heurística do protocolo e a inicialização no valor verdadeiro usada no
artigo.

\section{Métricas e testes}

São reportados MSE por parâmetro e agregado, erro absoluto mediano, viés, validade,
erro da razão de ramificação, custo e convergência do MLE. Como os parâmetros têm
escalas diferentes, também se reporta o \(R^2\) de cada parâmetro em relação à média
do prior (o estimador trivial). A comparação entre redes usa duas fontes de variação:
entre seeds (teste de Welch e de Mann--Whitney com cinco seeds), que mede a
variabilidade do treinamento \citep{bouthillier2021}, e entre trajetórias do teste
(diferença pareada com IC \emph{bootstrap} e teste de Wilcoxon), que mede a precisão
dos modelos já treinados \citep{efron1993,diebold1995}.

\chapter{Resultados}
\label{cap:resultados}

\section{Gate de dados}

Todos os 56 arquivos principais existiam. Foram geradas 28 partições. BTCUSDT produziu
887.803 eventos; ETHUSDT, 939.101. A janela independente gerou 1.288.352 eventos.
O gate de disponibilidade e volume foi, portanto, satisfeito.

\section{Reprodução no ponto fixo}

\input{e2_fixed_point.tex}

A Tabela~\ref{tab:e2-ponto-fixo} reproduz o padrão qualitativo de \citet{lee2023}.
O MLE é praticamente não viesado, e seus desvios-padrão são cerca de
\(\sqrt{2}\) vezes os do artigo, o esperado para trajetórias com metade do
comprimento. As redes têm o viés característico de encolhimento: aproximam
\(\alpha_{\mathrm{self}}\) e \(\alpha_{\mathrm{cross}}\) do centro do prior, como no
artigo, em troca de menor variância. As duas versões do MLE produziram exatamente as
mesmas estimativas, de modo que a inicialização no valor verdadeiro usada por Lee não
explica a vantagem do MLE nesse ponto bem identificado (\(g\approx0{,}73\)). O MSE
macro foi 0,0103 para o MLE e 0,0111 para o ensemble da LSTM.

\section{Comparação final: LSTM, Transformer e MLE}

\input{final_results.tex}

\input{posthoc_results.tex}

\section{Custo computacional}
\label{sec:custo}

A execução final registrou 105~ms de parede por trajetória para o MLE com quatro
processos. Esse valor foi medido enquanto o treino neural ocupava a GPU. Um benchmark
controlado na mesma máquina, com o mesmo código (que reproduz exatamente as 15.000
estimativas registradas), obteve os valores da Tabela~\ref{tab:custo}.

\begin{table}[htbp]
\centering
\caption{Custo por trajetória de 1.000 eventos em benchmark controlado.}
\label{tab:custo}
\small
\begin{tabular}{lr}
\toprule
Estimador e modo de execução & Tempo por trajetória\\
\midrule
MLE, 5 inícios, um núcleo & 9,6 ms\\
MLE, 1 início, um núcleo & 2,2 ms\\
MLE, 5 inícios, 4 processos (parede) & 2,5 ms\\
LSTM, lote unitário, CPU (1 thread) / MPS & 7,0 ms / 29,6 ms\\
Transformer, lote unitário, CPU (1 thread / 10 threads) / MPS & 9,9 / 4,8 / 5,0 ms\\
LSTM, lote de 256 no MPS & 0,07 ms\\
Transformer, lote de 32 no MPS & 3,7 ms\\
\midrule
\citet{lee2023}: rede / MLE & 12 ms / 1.763 ms\\
\bottomrule
\end{tabular}
\end{table}

Com gradiente analítico e recursão linear, o MLE desse modelo de quatro parâmetros
custa poucos milissegundos, a mesma ordem da inferência neural em lote unitário. A
razão de cerca de 150 vezes reportada por Lee não se repete. A vantagem de custo das
redes aparece na inferência em lote, sobretudo para a LSTM, e cresceria em modelos com
verossimilhança mais cara.

\section{Calibração empírica}

O MLE nas 1.063 janelas de calibração convergiu em todas. Os detalhes estão na
Seção~\ref{sec:calibracao}: a excitação cruzada é praticamente nula, \(\beta\) supera o
limite do prior na maior parte das janelas e nenhuma janela cai dentro de todas as
faixas usadas no treino.

\chapter{Discussão}

\section{Avaliação das hipóteses}

\textbf{H1 --- suportada quanto a estrutura e volume.} Os dados públicos permitem
reconstruir sequências com a mesma estrutura das de Lee, em volume suficiente. Seu
conteúdo estatístico, porém, difere do de ações: direção persistente, saltos de vários
ticks e tempos discretizados pela grade (Seção~\ref{sec:caracteristicas}).

\textbf{H2 --- parcialmente suportada.} O erro da LSTM é da mesma ordem do MLE: no
ponto fixo, 0,0111 contra 0,0103; no teste sintético, a perda mediana por trajetória
é praticamente igual. A latência substancialmente menor não se confirmou em lote
unitário frente a um MLE eficiente (Seção~\ref{sec:custo}); ela vale para inferência
em lote.

\textbf{H3 --- não suportada no sintético.} O Transformer não reduziu o erro: o MSE
macro médio foi maior que o da LSTM, sem significância entre seeds, e o ensemble
Transformer perdeu para o ensemble LSTM em todos os regimes. Ele também não se saiu
melhor no regime quase crítico, em que a memória longa seria mais relevante. Sua
vantagem foi a latência no MPS (5 ms contra 30 ms da LSTM), ao custo de treino
2,7 vezes maior. O resultado vale para a configuração testada; não houve busca de
hiperparâmetros, e o Transformer \emph{post-LN} foi treinado sem aquecimento da taxa
de aprendizado.

\section{Por que as redes vencem no MSE agregado}

O MSE macro do MLE (0,506) é cerca de 25 vezes o das redes, mas essa diferença
decorre de dois mecanismos, e não de maior precisão das redes na trajetória típica.
O primeiro é a identificabilidade (Seção~\ref{sec:identificabilidade}): no regime de
\(g\) baixo, o MLE produz valores extremos de \(\beta\) que maximizam legitimamente a
verossimilhança, e 1\% das trajetórias concentra 94\% do erro quadrático de \(\beta\).
O segundo é o prior (Seção~\ref{sec:bayes}): as redes aproximam a média a posteriori
sob o prior de treino e nunca saem do intervalo \([0{,}8;\,2{,}0]\) de \(\beta\), o que
limita seu pior erro. Essa robustez é real e útil, mas é comprada com viés de
encolhimento e depende de o prior representar os dados.

\section{Comparação com Lee (2023)}

As conclusões qualitativas de Lee se mantêm: a rede atinge erro da mesma ordem do MLE,
com viés de encolhimento, e o MLE é quase não viesado quando o parâmetro é bem
identificado. Os valores absolutos de MSE não são comparáveis entre os trabalhos,
pois dependem das faixas de parâmetros, que o artigo não publica. Duas divergências
merecem registro: o ganho de velocidade de duas ordens de grandeza não se repete com
um MLE eficiente, e a avaliação apenas por MSE agregado esconde as falhas de
identificabilidade que a inicialização no valor verdadeiro tende a mascarar.

\section{Implicações para a etapa empírica}

Os parâmetros estimados em janelas reais (\(\alpha_{\mathrm{cross}}\approx0\),
\(\beta\approx2{,}7\)) ficam fora do suporte do prior de treino, e as redes atuais
encolheriam suas estimativas para dentro desse suporte. A excitação cruzada nula
contrasta com o ponto típico de ações, em que ela domina (no ponto fixo de Lee,
\(\alpha_{\mathrm{cross}}>\alpha_{\mathrm{self}}\)), e é coerente com a persistência
de direção medida nos dados. Antes de qualquer comparação empírica, portanto, é
preciso recalibrar o prior com as janelas de calibração, ou treinar com rótulos MLE
em janelas reais, como sugere \citet{lee2023}. A discretização em 100~ms indica uma
terceira via: simular e filtrar na mesma grade, caso em que o MLE de tempo contínuo
fica mal especificado e a estimação amortizada ganha vantagem estrutural
\citep{lambe2025}.

\section{Ameaças à validade}

\begin{itemize}
  \item \textbf{Interna:} vazamento por janelas sobrepostas; tuning após observar o
        teste. Mitigadas por splits cronológicos e teste fechado por flag.
  \item \textbf{Construto:} o MLE pode ser rótulo imperfeito sob má especificação; nas
        janelas reais, $\alpha_{\mathrm{cross}}$ fica na fronteira do espaço de
        parâmetros, em que a teoria assintótica usual não vale.
  \item \textbf{Prior:} as métricas sintéticas dependem das faixas sorteadas, que não
        foram calibradas nos dados.
  \item \textbf{Externa:} dois ativos e duas janelas não representam todo o mercado.
  \item \textbf{Estatística:} cinco seeds quantificam apenas parte da variabilidade de
        treino; a LSTM apresentou dois níveis de convergência entre seeds, e uma seed
        atingiu o limite de 300 épocas ainda melhorando.
  \item \textbf{Reprodução:} o artigo omite faixas, seeds e detalhes de treino.
  \item \textbf{Computacional:} latência depende de hardware, lote, aquecimento e
        concorrência entre processos.
\end{itemize}

\chapter{Trabalhos futuros}
\label{cap:futuro}

\begin{enumerate}
  \item treinar as redes com rótulos MLE em janelas reais de calibração, selecionar na
        validação, abrir o teste empírico uma única vez e avaliar outubro com
        checkpoints congelados;
  \item recalibrar o prior sintético pelas janelas de calibração e retreinar;
  \item diagnosticar o ajuste nos dados reais por resíduos reescalados
        \citep{ogata1988,brown2002} e comparar com Poisson e com variantes do modelo;
  \item simular e filtrar na grade de 100~ms, para estimar sob a mesma observação dos
        dados reais;
  \item calcular volatilidade Hawkes e compará-la à realized volatility;
  \item padronizar os alvos, buscar hiperparâmetros e aquecer a taxa de aprendizado
        do Transformer, sem reutilizar o teste sintético para ajuste;
  \item adequar o documento ao modelo de TG do ITA.
\end{enumerate}

\chapter{Conclusão}

O trabalho construiu uma pipeline reproduzível que vai dos dados públicos da Binance
até a comparação de três estimadores sob protocolo congelado, e reproduziu o padrão
qualitativo de \citet{lee2023} no ponto fixo publicado. No teste sintético, a LSTM
não perdeu para o Transformer, e a atenção não trouxe ganho de precisão. A vantagem
das redes sobre o MLE no MSE agregado vem da robustez nas trajetórias em que
\(\beta\) é mal identificado e do encolhimento para o prior; na trajetória típica, o
MLE é igualmente preciso e, com implementação eficiente, igualmente rápido em lote
unitário. A etapa empírica mostrou que os dados de cripto ocupam outra região do
espaço de parâmetros, com excitação cruzada nula, o que torna a recalibração do prior
condição para aplicar as redes. A geração automática de tabelas e figuras reduz o
risco de divergência entre artefatos e texto.

\appendix
\chapter{Ambiguidades de Lee (2023)}

O texto menciona caminhos de 2.000 passos, mas uma forma de entrada com 1.000. Não
publica as faixas exatas de parâmetros, seeds, batch size nem dados NBBO. Essa menção
a 2.000 passos é mantida apenas como contexto histórico da reconstrução. O protocolo
final fixa 1.000 eventos e não compara comprimentos; decisões reconstruídas são
marcadas R2; cripto, A1; Transformer, E1. O artigo chama de ``três camadas'' as duas
camadas LSTM seguidas da camada densa; a implementação segue essa estrutura. A entrada
\(\log(1+\Delta t)\) e a codificação \(\pm1\) da direção também são R2.

\chapter{Reprodução computacional}

\begin{lstlisting}[language=bash]
uv sync --python 3.13 --extra dev
make test
make synthetic-final      # dados sinteticos com seed e SHA-256
make local-final          # treino retomavel, MLE, consolidacao, figuras e PDFs
make fixed-point          # experimento E2 no ponto fixo de Lee
make posthoc              # calibracao empirica e analise complementar
\end{lstlisting}

\chapter{Histórico de validação da pipeline}

Antes do protocolo congelado, testes reduzidos validaram a integração dos componentes.
Seus números não são comparáveis aos resultados finais e são registrados apenas para
rastreabilidade.

\begin{table}[htbp]
\centering
\caption{Testes de integração anteriores ao protocolo congelado.}
\label{tab:historico}
\small
\begin{tabularx}{\linewidth}{lXX}
\toprule
Teste & Configuração & Resultado\\
\midrule
Sintético reduzido & 1.024/128 trajetórias, 128 eventos, 3 épocas &
  MSE LSTM 0,0448; Transformer 0,0328\\
Integração empírica & 12 janelas rotuladas pelo MLE, só treino &
  MSE de treino 0,5038 e 0,1800\\
Protocolo congelado reduzido & 48/16/16 trajetórias, 2 épocas & pipeline validada\\
\bottomrule
\end{tabularx}
\end{table}

\bibliographystyle{plainnat}
\bibliography{references}
\end{document}
